\documentclass[12pt, twoside]{article}
\input{/Users/chris/github/course-files/docs/ib/preamble}

\fancyhead[LE]{\thepage}
\fancyhead[RO]{Markscheme}
\fancyhead[LO]{La Scuola d'Italia / Dr.\ Huson / IB Math 12 HL \\* 9 September 2026}

\begin{document}
\subsubsection*{12.2 Classwork: Complex Numbers Restart (no calculator)
--- Markscheme}

\begin{enumerate}[itemsep=0.3cm]

%--- Problem 1: Cartesian arithmetic, modulus, real unknowns [7 marks] ---
%--- Source: authored 2026-09-07 (HL syllabus 1.12) ---
\item[1.] [Maximum mark: 7]

$z = 3 - 2i$, $w = 1 + 4i$.

\medskip

\textbf{(a)} Find $zw$. \hfill [2]

\textbf{Solution:}

$(3 - 2i)(1 + 4i) = 3 + 12i - 2i - 8i^2$ \hfill \textbf{M1}

$zw = 11 + 10i$ \hfill \textbf{A1}

\emph{M1 for any complete expansion of the product, however the four terms
are organised.}

\medskip

\textbf{(b)} Find $\dfrac{z}{w}$ in the form $a + bi$. \hfill [2]

\textbf{Solution:}

$\dfrac{3 - 2i}{1 + 4i} \cdot \dfrac{1 - 4i}{1 - 4i}
= \dfrac{3 - 12i - 2i + 8i^2}{1 + 16}$ \hfill \textbf{M1}

$\dfrac{z}{w} = \dfrac{-5 - 14i}{17} = -\dfrac{5}{17} - \dfrac{14}{17}i$
\hfill \textbf{A1}

\emph{M1 for any correct method of realising the denominator (multiplying
by the conjugate, or equating $(a+bi)(1+4i) = 3-2i$).}

\emph{Accept $\dfrac{-5 - 14i}{17}$ or any equivalent exact form. Do not
accept a decimal approximation.}

\medskip

\textbf{(c)} Find the exact value of $|z|$. \hfill [1]

\textbf{Solution:}

$|z| = \sqrt{3^2 + (-2)^2} = \sqrt{13}$ \hfill \textbf{A1}

\medskip

\textbf{(d)} Find real $x$, $y$ with $(x + yi)(2 - i) = 7 + 4i$. \hfill [2]

\textbf{Solution:}

$(x + yi)(2 - i) = 2x + y + (2y - x)i$

Equate real and imaginary parts: $2x + y = 7$ and $2y - x = 4$
\hfill \textbf{M1}

$x = 2$, $y = 3$ \hfill \textbf{A1}

\emph{M1 for any method that separates the equation into two real
equations. A1 requires both values; the pair comes from one simultaneous
system, so the two marks are not split.}

\emph{Accept a solution obtained by dividing: $x + yi = \dfrac{7+4i}{2-i}$,
provided both values are correct.}

\newpage

%--- Problem 2: Non-real roots, conjugate pairs, sum and product [7 marks] ---
%--- Source: authored 2026-09-07 (HL syllabus 1.12) ---
\item[2.] [Maximum mark: 7]

$z^2 - 4z + 13 = 0$, $z \in \mathbb{C}$.

\medskip

\textbf{(a)} Solve the equation. \hfill [3]

\textbf{Solution:}

$(z - 2)^2 + 9 = 0$ \hfill \textbf{M1}

$(z-2)^2 = -9$

$z = 2 + 3i$, $z = 2 - 3i$ \hfill \textbf{A1A1}

\emph{M1 for completing the square or for the quadratic formula; either
route earns it. A1 for each root; accept $2 \pm 3i$ as a single
expression.}

\medskip

\textbf{(b)} Explain why the two roots are complex conjugates. \hfill [1]

\textbf{Solution:}

The coefficients of the equation are real, so any non-real root occurs
together with its conjugate. \hfill \textbf{R1}

\emph{R1 for any statement of the conjugate-root property, or for the
equivalent observation that the roots differ only in the sign of the
$\pm 3i$ term. A bare restatement of the two roots earns no mark.}

\medskip

\textbf{(c)} Sum and product of the roots of the equation with root
$2 + i\sqrt{3}$. \hfill [2]

\textbf{Solution:}

The other root is $2 - i\sqrt{3}$.

Sum $= (2 + i\sqrt{3}) + (2 - i\sqrt{3}) = 4$ \hfill \textbf{A1}

Product $= 2^2 - (i\sqrt{3})^2 = 4 + 3 = 7$ \hfill \textbf{A1}

\emph{A1 for each of the sum and the product; the two are independent. The
conjugate root must be used, stated or implied; a sum or product built from
any other second root earns no mark.}

\medskip

\textbf{(d)} Write down such a quadratic equation. \hfill [1]

\textbf{Solution:}

$z^2 - 4z + 7 = 0$ \hfill \textbf{A1}

\emph{Accept any non-zero real multiple, e.g.\ $2z^2 - 8z + 14 = 0$, and
accept the factored form $(z - 2 - i\sqrt{3})(z - 2 + i\sqrt{3}) = 0$.}

\emph{FT: award A1 for $z^2 - (\text{candidate's sum})z +
(\text{candidate's product}) = 0$ formed correctly from part (c).}

\end{enumerate}
\end{document}